This section presents visual evidence complementing the analysis from
Section~\ref{sec:discretization} and Section~\ref{sec:validation}.

Figure~\ref{fig:discretization-nodes} shows the node placement on the
uniform temporal mesh defined in Eq.~\eqref{eq:time-grid}. Increasing $N$
reduces the spacing $h$, which provides more sample points for resolving the
decay curve.

\begin{figure}[h]
\centering
\includegraphics[width=0.90\textwidth]{docs/latex/source/figures/discretization_nodes.png}
\caption{Uniform-grid node distributions for coarse and refined discretizations.}
\label{fig:discretization-nodes}
\end{figure}

Figure~\ref{fig:trajectory-convergence} overlays numerical trajectories for
increasing $N$ against the analytical solution from Eq.~\eqref{eq:exact-solution}.
The curves visually confirm convergence under refinement.

\begin{figure}[h]
\centering
\includegraphics[width=0.90\textwidth]{docs/latex/source/figures/trajectory_convergence.png}
\caption{Numerical trajectories converging to the analytical solution as $N$ increases.}
\label{fig:trajectory-convergence}
\end{figure}

Figure~\ref{fig:endpoint-convergence} reports endpoint behavior for multiple
initial temperatures $T_0$. Each panel shows both endpoint approximation and
absolute endpoint error versus $N$, demonstrating convergence by refinement for
different initial states.

\begin{figure}[h]
\centering
\includegraphics[width=0.92\textwidth]{docs/latex/source/figures/endpoint_convergence_by_initial.png}
\caption{Endpoint convergence for multiple initial temperatures and grid sizes.}
\label{fig:endpoint-convergence}
\end{figure}
